\begin{abstract} % 150 words
% what ius this about
%what we did 
%how we did it
%why we did it
%what we found?
{\color{blue}
Reliable verification of diffusion-generated watermarks must remain effective under diverse image transformations, (i.e. compression, cropping, resampling, geometric distortions, and compound platform processing). Although existing watermarking methods can provide robust provenance signals, their performance may vary substantially across different transformations, making the construction of effective training regimes challenging.
In this paper, we introduce AI-Meta-Mark that treats different image transformations as distinct meta-training tasks and adaptively learns how these tasks should be sampled during verifier training. Our framework comprises two stages: $(1)$ algorithmic frequency-domain watermark injection, and $(2)$ learning-based probabilistic verifier within a meta-learning formulation. Specifically, in $(1)$,  a structured watermark is injected into the initial diffusion noise and propagated through the generation process, while in $(2)$, the verifier learns to recover the corresponding spectral signature. To coordinate meta-training tasks, we propose an Agentic Task Sampler that constructs attack-conditioned support and query sets and dynamically adjusts task sampling according to the following key factors: $(i)$ query loss, $(ii)$ verification failure rate, $(iii)$ uncertainty, and $(iv)$ sampling history. This closed-loop strategy enables the training process to focus on difficult, underrepresented, and informative transformations rather than relying on a fixed task distribution.
Extensive experiments on a 1,000-image COCO test set demonstrate that AI-Meta-Mark outperforms six alternative task-scheduling strategies across eight image transformations [results to be completed]. These results highlight the potential of agentic, attack-aware task scheduling for improving the robustness and generalization of diffusion-watermark verification.}

%During generation, a structured watermark is injected into the initial diffusion noise and propagated through the sampling process. During verification, a lightweight neural verifier recovers the corresponding spectral signature and produces a probabilistic confidence score.




\sout{
SpiderMark combines structured frequency-domain watermarking with a learned verifier, but fixed perturbation mixtures cannot adapt to changing attack difficulty. We propose \emph{Agentic MetaSpiderMark}, which uses online verifier feedback to schedule attack-conditioned meta-learning episodes while leaving the watermark and diffusion model unchanged. Because deployed verification requires a binary decision, we prioritize fixed-threshold accuracy and report AUROC as a secondary metric. In a matched comparison of six schedulers on an unseen 1,000-image COCO test set, Agentic MetaSpiderMark achieves the highest mean accuracy across eight attacks: 89.15$\%$ (95$\%$ bootstrap CI: 88.06--90.21). Paired bootstrapping confirms positive gains over uniform sampling, Adaptive Sampler, and Bandit-UCB. BASS and ATS have lower mean accuracy but remain statistically competitive. Uniform slightly leads AUROC (96.00$\%$ versus 95.86$\%$). Agentic MetaSpiderMark is therefore the strongest accuracy-first scheduler in this controlled benchmark, rather than a universal winner under every metric.}

%Reliable watermark verification for diffusion-generated images must remain stable after the image has passed through unknown downstream transformations, including compression, cropping, resampling, geometric warping, and compound platform-style processing.
%SpiderMark addresses this problem with structured frequency-domain watermark injection and a learned verifier, but the verifier is still trained from a fixed distribution of perturbations and checkpoints.
%This paper introduces \emph{Agentic MetaSpiderMark}, an agentic AI task scheduler for meta-learning-based watermark verification.
%The proposed scheduler monitors online verifier-training signals, including attack-specific query loss, loss change, uncertainty, failure rate, and sampling history, and converts them into task-weight updates for the next attack-conditioned episodes.
%Rather than using a fixed augmentation mixture or a hand-designed curriculum, Agentic MetaSpiderMark closes the loop between verifier feedback and downstream attack exposure while keeping the SpiderMark watermark pattern and diffusion generator unchanged.
%The framework constructs watermark verification episodes from attack-conditioned support and query sets, performs inner verifier adaptation, and lets the agentic scheduler decide which attack tasks should receive more training attention.
%We define a scheduler benchmark that compares the agentic scheduler against uniform sampling, bandit scheduling, ATS-style scheduling, BASS-style scheduling, and Adaptive Sampler-style task sampling under matched verifier episodes.
%Preliminary experiments on nine image transformations show that the meta-learned verifier improves mean accuracy from 0.8499 to 0.9182 and mean AUROC from 0.9293 to 0.9686 over a non-meta SpiderMark verifier, with gains on every evaluated attack.
%The largest observed improvement occurs under compound messaging-app processing, where accuracy rises from 0.6813 to 0.9173.
%These results suggest that robust watermark verification is not only a matter of watermark geometry or verifier architecture, but also of how an adaptive agent schedules downstream verification tasks during meta-training.
\end{abstract}
